\documentclass[11pt]{article}

\usepackage{cmap}
\usepackage[T1]{fontenc}
\usepackage{lmodern}
\input{glyphtounicode}
\pdfgentounicode=1
\usepackage[margin=1in]{geometry}
\usepackage{amsmath,amssymb}
\usepackage{array}
\usepackage{booktabs}
\usepackage{xcolor}
\usepackage{upquote}
\usepackage{listings}
\usepackage[hidelinks]{hyperref}
\hypersetup{
  pdftitle={nsm: explainable menu-relative anomaly repair for U(1)X extensions of the Standard Model},
  pdfauthor={Evan Morritse}
}

\definecolor{codebg}{gray}{0.96}
\lstset{
  basicstyle=\ttfamily\small,
  backgroundcolor=\color{codebg},
  breaklines=true,
  frame=none,
  columns=fullflexible,
  keepspaces=true,
  xleftmargin=1em,
}

\newcommand{\nsm}{\texttt{nsm}}
\newcommand{\UY}{U(1)_Y}
\newcommand{\UX}{U(1)_X}
\newcommand{\SUthree}{SU(3)}
\newcommand{\SUtwo}{SU(2)}

\title{\nsm: explainable menu-relative anomaly repair\\
for $U(1)_X$ extensions of the Standard Model}
\author{Evan Morritse}
\date{\today}

\begin{document}
\maketitle

\begin{abstract}
This paper presents \nsm, an explainable symbolic engine for diagnosing and repairing
anomalous $U(1)_X$ extensions of the Standard Model. Given a one-generation
charge assignment, \nsm{} evaluates the complete
$SU(3)\times SU(2)\times U(1)_Y\times U(1)_X$ perturbative local anomaly set,
including the mixed $[Y^2X]$ and $[YX^2]$ coefficients, identifies the
obstructing anomaly coefficients, and searches a declared menu of sterile and
SM-vector-like matter in increasing declared cost. The repair tier is
bounded below by exact support criteria: sterile singlets cannot cancel mixed hypercharge
anomalies, colorless additions cannot cancel $[\SUthree^2X]$, and weak singlets
cannot cancel $[\SUtwo^2X]$. In the sterile tier, one- and two-sterile-singlet
repairs have closed-form rationality criteria, so the solver is needed only where these exact
checks run out. \nsm{} returns the first exact rational repair found and certifies
its minimum cost when every cheaper rung is resolved by support pruning or SMT
unsatisfiability. Every returned repair spectrum is re-verified against all
anomaly coefficients, while unresolved solver cases are marked uncertified. The
minimality claim is relative to the declared N/E'/L'/D'/Q' representation menu,
the stated lexicographic cost ordering, and the rational-charge requirement. The
method is validated by deriving Standard-Model hypercharges, reproducing the
one-$\nu_R$-per-generation repair of family-universal gauged $B\!-\!L$, and
constructing a repair census over bounded integer charge assignments. The method
intentionally does not synthesize scalar sectors, enforce $\UX$ Yukawa
completion for repaired models, or calculate physical observables.
\end{abstract}

\noindent\textbf{Keywords:} Anomaly cancellation; Gauge anomaly; \(U(1)\) gauge extension; Symbolic computation; Satisfiability modulo theories; Computational particle physics

\section*{Program Summary}

\noindent
\textbf{Program title:} \texttt{nsm}

\noindent
\textbf{CPC Library link to program files:} To be assigned by CPC if accepted.

\noindent
\textbf{Developer's repository link:}
\url{https://github.com/Morritse/nsm-anomaly-repair}.

\noindent
\textbf{Code Ocean capsule:} To be added by Technical Editor.

\noindent
\textbf{Licensing provisions:} MIT License.

\noindent
\textbf{Programming language:} Python 3.

\noindent
\textbf{External routines/libraries:} \texttt{z3-solver} for SMT solving.
Exact rational arithmetic uses the Python standard-library \texttt{fractions}
module. Development and validation use pytest.

\noindent
\textbf{Nature of problem:} For a proposed \(U(1)_X\) extension of the Standard
Model, determine whether the fermion spectrum is locally anomaly-free. If it is
not, search a declared menu of sterile and SM-vector-like but
\(U(1)_X\)-chiral matter for an exact rational anomaly-cancelling repair, while
recording which lower-cost repairs are ruled out by support arguments or solver
unsatisfiability.

\noindent
\textbf{Solution method:} The program evaluates the full perturbative local
anomaly vector in exact rational arithmetic, applies representation-theoretic
support bounds to prune impossible repair tiers, uses closed-form rationality
criteria for one- and two-sterile-singlet repairs, and sends remaining nonlinear
repair systems to SMT solving. Returned rational repair witnesses are
re-verified against all anomaly coefficients. Minimality is certified only when
every lower-cost rung is resolved by support pruning or solver unsatisfiability.

\noindent
\textbf{Additional comments including restrictions and unusual features:}
Minimality is relative to the declared
\(N/E'/L'/D'/Q'\) matter menu, the stated lexicographic cost ordering, the Weyl
multiplet budget, and the rational-charge requirement. The program does not
synthesize scalar sectors, enforce Yukawa completion for repaired models, or
compute physical observables. The repair result includes an evidence record
distinguishing certified minimality from unresolved solver rungs and algebraic
non-rational solver models.

\section{Introduction}

Most computational tools for particle theory are \emph{calculators} (amplitudes,
cross sections, decay rates) or \emph{databases} (particle properties, model
files). Comparatively little tooling answers the earlier structural questions a
model builder actually asks first: \emph{Given this gauge structure, what is
forced? What is merely conventional? Is this proposed extension even
consistent, and if not, what is the smallest thing that fixes it?}

\nsm{}\footnote{\nsm{} is the package name, not an acronym; the project began
as a ``non-standard model'' pun and is now used here as a consistency engine for
Standard-Model and BSM gauge theories.} is an attempt to make those questions
mechanical and explainable for the anomaly-repair problem. It combines exact
rational anomaly arithmetic, simple representation-theoretic support bounds, and
SMT search over the remaining nonlinear systems. The design goal is simple: when
the tool says that a charge assignment is broken, it should also say where it
broke, which cheaper repairs cannot reach the obstruction, and what kind of
repair can actually fix it. The output is therefore a \emph{traceable structural
statement}, rather than a floating-point number whose provenance is opaque.

\paragraph{Contributions.}
\begin{enumerate}
\item Exact anomaly diagnosis for
$SU(3)\times SU(2)\times U(1)_Y\times U(1)_X$, including the mixed $[Y^2X]$ and
$[YX^2]$ coefficients.
\item Support-based lower bounds on the representation tier required to repair a
residual anomaly vector, making many obstruction claims solver-free.
\item A matter-repair search over a declared menu of sterile
singlets and SM-vector-like, $\UX$-chiral exotic pairs, with closed-form sterile
criteria, exact rational witnesses, and certified minimum-cost records when all
cheaper rungs are resolved.
\item Worked repairs and a reproducible repair census, validated by exact
re-checking of every returned spectrum and by reproduction of known results.
\end{enumerate}

\paragraph{Non-contributions.}
\nsm{} computes no physical observables and makes no claim of discovery. It does
not derive masses, mixing angles, coupling values, scalar sectors, or the number
of generations (these remain empirical inputs). Its results are structural and
quantum-number-level; see Section~\ref{sec:limits}. Throughout, this paper uses the
convention $Q = T_3 + Y/2$, so the quark doublet carries $Y = 1/3$.

\section{Repair problem}

The input to the central algorithm is a one-generation $\UX$ charge assignment
for the five Standard-Model multiplets $(Q_L,u_R,d_R,L_L,e_R)$. Missing entries
are interpreted as zero charge, not as deleted fields. The task is to decide
whether the assignment is anomaly-free; if it is not, search a declared menu of
new Weyl multiplets for an anomaly-cancelling matter repair.

Equivalently, for an input spectrum $S$, declared repair menu $\mathcal M$,
local anomaly vector $\mathcal A$, and lexicographic cost $C$ defined in
Section~4, the repair problem is
\[
\underset{R\in\mathcal M}{\operatorname{minimize}}\; C(R)
\qquad\text{subject to}\qquad
\mathcal A(S\cup R)=0,\qquad X_R\in\mathbb Q .
\]
The output is not just a spectrum, but a rational witness $R$ together with its
certification status under that declared ordering.

The rationality requirement is physical as well as practical. For a compact
$\UX$ gauge factor, charges live on a lattice; after an overall normalization,
rational charges can be represented as integer lattice charges. Algebraic
irrational solver models are therefore treated as evidence of real
satisfiability, but not as accepted compact-$\UX$ repair witnesses.

The menu used in this paper is shown in Table~\ref{tab:menu}. $N$ is a sterile
right-handed singlet. The other entries are SM-vector-like pairs: the left and
right Weyl multiplets carry the same SM representation but may carry different
$\UX$ charges, so they are chiral under the new gauge factor while remaining
vector-like under the Standard Model. Throughout, ``six added fields'' means six
added \emph{Weyl multiplets}, not six individual color or weak components.

\begin{center}
\centering
\small
\begin{tabular}{@{}l>{\raggedright\arraybackslash}p{0.24\linewidth}c>{\raggedright\arraybackslash}p{0.40\linewidth}@{}}
\toprule
item & SM representation(s) & Weyl multiplets & anomaly support \\
\midrule
$N$ & $(1,1,0)_R$ & $1$ & $[\text{grav}^2X]$, $[X^3]$ \\
$E'$ & $(1,1,-2)_L+(1,1,-2)_R$ & $2$ &
$[\text{grav}^2X]$, $[Y^2X]$, $[YX^2]$, $[X^3]$ \\
$L'$ & $(1,2,-1)_L+(1,2,-1)_R$ & $2$ &
$[\SUtwo^2X]$, $[\text{grav}^2X]$, $[Y^2X]$, $[YX^2]$, $[X^3]$ \\
$D'$ & $(3,1,-2/3)_L+(3,1,-2/3)_R$ & $2$ &
$[\SUthree^2X]$, $[\text{grav}^2X]$, $[Y^2X]$, $[YX^2]$, $[X^3]$ \\
$Q'$ & $(3,2,1/3)_L+(3,2,1/3)_R$ & $2$ &
$[\SUthree^2X]$, $[\SUtwo^2X]$, $[\text{grav}^2X]$, $[Y^2X]$, $[YX^2]$, $[X^3]$ \\
\bottomrule
\end{tabular}
\par\medskip
\refstepcounter{table}\label{tab:menu}
Table~\thetable: Declared repair menu. The support column lists which residual
$X$-sector anomaly coefficients an item can change with a chiral $\UX$ split.
\end{center}

\section{Anomaly conventions}

For a newly \emph{gauged} $\UX$ alongside hypercharge, the engine evaluates the
complete perturbative local anomaly set for
$\SUthree\times\SUtwo\times\UY\times\UX$, together with the Witten $\SUtwo$
doublet-parity check~\cite{witten}. The local set includes pure-$X$ and
\emph{mixed} coefficients. Writing $d_3,d_2$ for the color and weak dimensions,
$T_c,T_w$ for the Dynkin indices, and $s=\pm1$ for chirality, the engine
evaluates
\begin{align}
[\SUthree^2 X] &= \textstyle\sum s\,d_2 T_c\,X, &
[\SUtwo^2 X] &= \textstyle\sum s\,d_3 T_w\,X, &
[\text{grav}^2 X] &= \textstyle\sum s\,d_3 d_2\,X, \nonumber\\
[X^3] &= \textstyle\sum s\,d_3 d_2\,X^3, &
[Y^2 X] &= \textstyle\sum s\,d_3 d_2\,Y^2 X, &
[Y X^2] &= \textstyle\sum s\,d_3 d_2\,Y X^2,
\label{eq:mixed}
\end{align}
together with the corresponding $Y$-sector and $[\SUthree^3]$ coefficients. The
mixed terms~\eqref{eq:mixed} are load-bearing: a charge assignment can cancel
every pure-$X$ anomaly yet fail $[Y^2 X]$, so a single-$U(1)$ check can certify
a still-anomalous model as consistent.

\section{Minimal repair}

Given a charge assignment that fails the conditions above, \nsm{} searches for
the smallest set of added fermions whose charges cancel the residual anomalies.
Some of this reasoning is solver-free. For each menu item define its
\emph{support}: the set of anomaly coefficients it can change with a chiral
$\UX$ split. A sterile singlet has support
\[
\{[\text{grav}^2X], [X^3]\},
\]
because $Y=0$ and it is an $\SUthree\times\SUtwo$ singlet. A colorless
hypercharged pair reaches the mixed hypercharge coefficients but not
$[\SUthree^2X]$; a weak singlet cannot reach $[\SUtwo^2X]$; and a colored pair
is required to change $[\SUthree^2X]$. Therefore, if the support of the residual
anomaly vector is not contained in the union of the chosen menu supports, no
charge assignment in that menu can repair the model. This gives exact lower
bounds on the repair tier before any SMT query is made.

The sterile tier also has a closed form. A sterile right-handed singlet is color
and weak neutral with $Y=0$ and unknown $X$, so cancelling a residual
$(g,c)=([\text{grav}^2X],[X^3])$ requires
\begin{equation}
\sum_i X_i = g, \qquad \sum_i X_i^3 = c.
\label{eq:repair}
\end{equation}
One sterile field exists iff $c=g^3$ (with the field charge $X=g$). For two
sterile charges, write $x+y=g$ and $x^3+y^3=c$. If $g\ne0$, then
\[
xy=\frac{g^3-c}{3g}, \qquad
\Delta=(x-y)^2 = g^2-4xy = \frac{4c-g^3}{3g}.
\]
A rational two-singlet repair exists iff $\Delta$ is a rational square (with
the usual nonzero-field guard if inert zero-charge additions are disallowed).
If $g=0$, then $x=-y$ forces $c=0$, so a two-field sterile pair cannot cancel a
nonzero cubic residual. Thus the first two sterile rungs are decidable by hand;
higher-multiplicity sterile systems and exotic-pair systems are passed to the
solver.

Beyond the singlet, the menu adds \emph{SM-vector-like, $\UX$-chiral} exotic
pairs $E',L',D',Q'$ in definite SM representations with left/right charges
$X_L\neq X_R$. Being vector-like under the Standard Model, each pair cancels its
own SM anomalies; its chiral $\UX$ split feeds the $X$-sector, with contributions
$\propto(X_L-X_R)$, $(X_L^2-X_R^2)$, and $(X_L^3-X_R^3)$. Hypercharged pairs
therefore reach $[Y^2X]$ and $[YX^2]$; colored pairs reach $[\SUthree^2X]$; weak
doublets reach $[\SUtwo^2X]$. The chirality condition $X_L\neq X_R$ is
essential: an $X$-vector-like pair would be inert in every $X$ anomaly.

The engine searches repair tiers in order---sterile singlet, then colorless
exotic, then colored exotic---and reports the first exact rational repair found
under the declared cost ordering. If a nonzero anomaly lies outside the reach of
every menu item, the model is reported \emph{structurally blocked} with that
coefficient named; if every nonzero anomaly is reachable but no repair fits the
Weyl-multiplet budget, it is reported \emph{budget-limited} instead. The tool
thus never claims a fix it cannot deliver, nor mistakes a budget limit for an
obstruction.

The ordering is lexicographic. First compare repair tier,
\[
N \;<\; \{N,E',L'\} \;<\; \{N,E',L',D',Q'\}.
\]
Within a tier, compare
\[
C(R)=\bigl(N_{\rm Weyl},\ \max_i w_i,\ \sum_i n_i w_i,\ \text{combination}\bigr),
\]
where $N_{\rm Weyl}=n_N+2(n_E+n_L+n_D+n_Q)$,
\[
\text{combination}=(n_N,n_E,n_L,n_D,n_Q),
\]
and
\[
w_N=1,\quad w_E=2,\quad w_L=4,\quad w_D=6,\quad w_Q=12 .
\]
The weights count the number of SM components carried by one menu item: one
sterile Weyl multiplet for $N$, two singlets for $E$, two weak doublets for $L$,
two color triplets for $D$, and two color-triplet weak-doublets for $Q$. The
final combination component is only a deterministic tie-breaker.

\nsm{} searches candidate spectra in increasing declared cost, returns the first
exact rational repair found, and certifies its minimum cost when every cheaper
rung is resolved by support pruning or SMT unsatisfiability. Where SMT is used,
the search emits a \emph{minimality evidence record}: for
each lower-cost menu combination under the declared ordering, it records whether
the combination was ruled out by support (\texttt{support-prune}), ruled out by an SMT \texttt{unsat}
verdict, left unresolved (\texttt{unknown}), or produced only an algebraic model
from the solver. An SMT \texttt{unsat} verdict is a decision over the reals, so a
record with only \texttt{support-prune} and \texttt{unsat} lower rungs is a
solver-backed guarantee that no cheaper repair exists within the menu. An
\texttt{unknown} or \texttt{sat-algebraic-model} rung does not certify
minimality and is reported as such. The implementation records solver verdicts, but not tracked
unsat cores, SMT-LIB replay artifacts, or proof logs; an unsat core would
identify a conflicting subset of assumptions, while a proof log, where supported,
would provide the stronger artifact. For gauged $B\!-\!L$ the criterion is
immediate: $g=c=-1$ and $c=g^3$, so the minimal one-field sterile repair is
a single $\nu_R$ with $X=-1$.

The reported repair is an anomaly-cancelling rational witness within the menu;
it is a certified minimum only when the evidence record has no unresolved
lower-cost rung. It is not unique, not a naturalness statement, and not a claim
that the field exists.

\section{\texorpdfstring{Case study: gauging $B\!-\!L$}{Case study: gauging B-L}}
\label{sec:casestudy}

Assign each fermion its $B\!-\!L$ charge ($+1/3$ for quarks, $-1$ for leptons)
and gauge it. The computed local anomaly set (Table~\ref{tab:bml}) shows the model is
inconsistent, and \emph{why}: only $[\text{grav}^2 X]$ and $[X^3]$ fail, both
equal to $-1$; every mixed and non-abelian coefficient already cancels.

\begin{table}[t]
\centering
\small
\begin{tabular}{@{}lrlr@{}}
\toprule
coefficient & value & coefficient & value\\
\midrule
$[\SUthree^3]$ & $0$ & $[\text{grav}^2 X]$ & $-1$ \ \textbf{(fails)}\\
$[\SUthree^2 Y]$, $[\SUthree^2 X]$ & $0$ & $[Y^3]$ & $0$\\
$[\SUtwo^2 Y]$, $[\SUtwo^2 X]$ & $0$ & $[Y^2 X]$, $[Y X^2]$ & $0$\\
$[\text{grav}^2 Y]$ & $0$ & $[X^3]$ & $-1$ \ \textbf{(fails)}\\
\bottomrule
\end{tabular}
\caption{The eleven anomaly coefficients for gauged $B\!-\!L$ with
Standard-Model fermions only, as computed by \nsm. Only the gravitational and
cubic $X$-anomalies fail.}
\label{tab:bml}
\end{table}

By~\eqref{eq:repair} with $g=c=-1$, the one-$\nu_R$-per-generation sterile repair
has $X=-1$: one sterile right-handed singlet of $B\!-\!L$ charge $-1$. Re-running
the local anomaly set on the repaired fermion content returns all zeros. The
familiar right-handed neutrino appears here as the answer to a constrained repair
problem, not as an input assumption.
Constructing a complete massive $B\!-\!L$ gauge model additionally requires a
symmetry-breaking scalar sector, which lies outside the present anomaly-repair
problem.

Gauged $B\!-\!L$ is a clean corner of the repair map: the obstruction is narrow,
and the minimal sterile repair is the textbook one~\cite{bminusL}. The same
engine is more useful when the obstruction is less tidy. It reports the
repair tier found for any $\UX$ charge assignment
(Table~\ref{tab:tiers}), and a second worked case makes the tiering
concrete. Assign $X=1$ to $e_R$ alone; the engine returns the anomaly vector
\[
\bigl([\text{grav}^2X],\,[Y^2X],\,[YX^2],\,[X^3]\bigr) = (-1,\,-4,\,+2,\,-1),
\qquad [\SUthree^2X]=[\SUtwo^2X]=0 .
\]

A sterile singlet has $Y=0$, so it feeds only $[\text{grav}^2X]$ and $[X^3]$ and
\emph{cannot} touch the mixed $[Y^2X]$ or $[YX^2]$; the engine reports the
sterile menu as structurally blocked on exactly those two coefficients. The
search is therefore forced up one tier, and the colorless menu repairs the model
with one hypercharged $E'_L/E'_R$ pair ($Y=-2$): two Weyl multiplets carrying
$X_L=1$, $X_R=0$ in the signed left/right convention of Eq.~\eqref{eq:mixed},
vector-like under the Standard Model and chiral under $\UX$. With that pair
included, all eleven coefficients re-verify to zero.

This example isolates anomaly repair. Invariance of the ordinary electron
Yukawa interaction under the added $\UX$ is a separate model-building constraint
unless the scalar charges or scalar sector are included in the skeleton.

The sterile tier is ruled out analytically: no number of sterile multiplets can
generate the $[Y^2X]$ or $[YX^2]$ support it must cancel. The two-field
colorless pair is therefore certified minimal within the declared menu and cost
ordering because all cheaper tiers are support-pruned; re-verification proves
the returned witness itself cancels all anomalies.

As a colored example, assign $X=1$ to $Q_L$ alone. The computed residual vector
has
\[
\begin{aligned}
\relax[\SUthree^2X]&=1, &
[\SUtwo^2X]&=\tfrac32, &
[\text{grav}^2X]&=6,\\
[Y^2X]&=\tfrac23, &
[YX^2]&=2, &
[X^3]&=6 .
\end{aligned}
\]
The nonzero $[\SUthree^2X]$ coefficient rules out the sterile and colorless
tiers by support before any nonlinear solve is attempted. A single $D'$ pair has
color support but not weak support, so it is also support-pruned at the same
field count. The first rational repair in the declared ordering is a colored
weak-doublet $Q'_L/Q'_R$ pair with $Y=1/3$ and $(X_L,X_R)=(0,1)$, after which
all eleven anomaly coefficients re-verify to zero.

\begin{table}[t]
\centering
\small
\begin{tabular}{@{}lll@{}}
\toprule
charge assignment & reported repair tier & example repair\\
\midrule
$B\!-\!L$ & sterile-repairable & $\nu_R$ with $X=-1$\\
$X$ on $e_R$ only & colorless-exotic-repairable & an $E'_L/E'_R$ chiral split\\
$X$ on $Q_L$ only & colored-exotic-required & a $Q'_L/Q'_R$ chiral split\\
\bottomrule
\end{tabular}
\caption{The reported repair tier for three $\UX$ charge assignments, as
classified by \nsm{}. $B\!-\!L$ is the familiar sterile case; the others demand
hypercharged or colored exotics that the sterile menu cannot supply.}
\label{tab:tiers}
\end{table}

This is not a particle-discovery machine. It is a consistency and repair map
within a declared representation menu and cost ordering. That turns the layer
from a checker that validates one known result into a tool that maps a space of
anomalous extensions into repair classes, reporting what \emph{kind} of new
matter a broken gauge extension demands when the lower-cost record is certified.

\paragraph{A repair census.}
Because the classifier is mechanical, it can be run over a family of extensions
rather than a single example. The census enumerates every integer $\UX$ charge
assignment to one generation with $|q|\le 2$, reduces rescaling/sign duplicates
to the zero assignment plus $1441$ nonzero primitive sign-fixed representatives
(the tier is invariant under $X\to\lambda X$; it is also invariant under
hypercharge mixing $X\to X+\alpha Y$, which this paper does not quotient but
uses as an internal consistency check), and classifies each representative by
its reported repair class. The
resulting census (Table~\ref{tab:census}) is a reproducible, basis-dependent
slice: one cell is the trivial zero generator, $0.1\%$ are sterile-repairable,
$9.6\%$ need only colorless exotics, $31.5\%$ require colored exotics, and
$58.7\%$ have no repair found with at most six added Weyl multiplets.

\begin{table}[t]
\centering
\small
\begin{tabular}{@{}lrrrr@{}}
\toprule
reported class & total & share & certified min./search & uncertified rung\\
\midrule
already-consistent & $1$ & $0.1\%$ & $1$ & $0$\\
sterile-repairable & $1$ & $0.1\%$ & $1$ & $0$\\
colorless-exotic-repairable & $139$ & $9.6\%$ & $139$ & $0$\\
colored-exotic-required & $454$ & $31.5\%$ & $454$ & $0$\\
no repair found ($\le6$ multiplets) & $847$ & $58.7\%$ & $766$ & $81$\\
\bottomrule
\end{tabular}
\caption{Repair census of the zero assignment plus $1441$ nonzero
primitive sign-fixed integer $\UX$ assignments to one generation with $|q|\le2$
in the fixed $(Y,X)$ basis (menu N/E/L/D/Q, budget six added Weyl multiplets,
solver timeout $8\,$s per SMT call). The raw grid has $3125$ cells ($2.17\times$
deduplication). ``Certified min./search'' means no lower-cost support or solver
rung was unresolved; for the last row it certifies only the bounded search
record, not existence or nonexistence above the budget.}
\label{tab:census}
\end{table}

Three caveats are built into the table. First, the last row is a
\emph{Weyl-multiplet-budget artifact}, not an unrepairability theorem: with the full menu
every failing anomaly is individually reachable (a necessary, not sufficient,
condition for a joint repair), no cell is structurally blocked, and the
unrepaired-within-budget fraction decreases in lower-budget runs as the budget
grows. The census therefore reports ``no repair found with
$\le6$ multiplets,'' not a proof that a repair exists or does not exist at larger
size. Second, the percentages count a bounded integer box in one fixed $(Y,X)$
basis; they are not a measure over all $\UX$ extensions, and a different choice
of generator would re-slice the population. Third, $81$ cells carry an
uncertified solver rung (an \texttt{unknown} or \texttt{sat-algebraic-model}
verdict in the minimality record), so their reported class is provisional. With
those stated, the census is the method claim made concrete: not one rediscovery,
but a reproducible map of which kind of new matter a bounded space of anomalous
gauge extensions demands when the record is certified, or appears to demand when
it is not.

\section{Implementation}
\label{sec:impl}

\nsm{} is implemented in Python over exact rational arithmetic, Z3, and small
graph searches; the narrated examples run via \texttt{python -m nsm}. In the
primary anomaly/repair path, consistency conditions are compiled to solver
assertions, and the routine that \emph{checks} a concrete spectrum is the same
one that \emph{solves} for an unknown one: the anomaly-coefficient arithmetic is
polymorphic in its charge arguments, returning exact rationals when fed
rationals (a check) and symbolic expressions when fed solver variables (a
constraint). Listing~\ref{lst:api} shows a representative query.

\begin{lstlisting}[language=Python,caption={Public API: forced hypercharges, and a minimal repair.},label=lst:api]
from fractions import Fraction as F
from nsm import derive_hypercharges, minimal_repair
from nsm.sm import standard_model_skeleton

r = derive_hypercharges(standard_model_skeleton())
r.hypercharges["Q_L"], r.is_unique     # -> (Fraction(1, 3), True)

# gauged B-L (charges as exact Fractions); repair searches for a fix
B_minus_L = (
    ("Q_L", F(1, 3)), ("u_R", F(1, 3)), ("d_R", F(1, 3)),
    ("L_L", F(-1)), ("e_R", F(-1)),
)
minimal_repair(B_minus_L).added        # -> [('nu_R', Fraction(-1, 1))]
\end{lstlisting}

The repair-census layer uses the same primitive over an integer charge grid: it
reduces each vector to a primitive sign-fixed representative, keeps the raw and
deduplicated denominators, records the representative's reported repair class,
and exports CSV, Markdown, or \LaTeX{} summaries. Hypercharge-basis shifts
are treated as an equivalence check rather than as a counting quotient, since
they shear the integer lattice.

The kernel poses the hypercharge problem as a conjunction of assertions (each
anomaly coefficient equal to zero, gauge invariance of each Yukawa coupling, and
one normalization) and reads the unique assignment from the model. Whether that
assignment \emph{is} unique is decided by a second solve: assert the negation of
the found solution and check satisfiability. \texttt{unsat} certifies
uniqueness; \texttt{sat} exhibits a distinct second solution (such as the
$u_R\!\leftrightarrow\!d_R$ ambiguity in the anomaly-only validation check); and
\texttt{unknown} (the solver exceeding its budget on the nonlinear real
arithmetic) is reported as a genuine third outcome rather than silently
collapsed to false certainty. The downstream structural routines use separate
combinatorial searches over configured vertices. The minimal-repair search
(Listing~\ref{lst:repair}) encodes anomaly cancellation for each candidate menu
combination, walks the declared cost ordering, and treats algebraic solver
models as uncertified for rational-repair purposes rather than failing on them.

\begin{lstlisting}[language=Python,caption={Tiered repair: search tiers, then the declared cost ordering within each tier.},label=lst:repair]
prior_records = []
for tier in [("sterile", N), ("colorless", NEL), ("colored", NELDQ)]:
    for cost, combo in get_combinations(max_added, tier.menu):
        if combo_cannot_reach_a_nonzero_anomaly(combo):
            prior_records.append(record(tier, cost, combo, "support-prune"))
            continue
        status, charges = solve_anomaly_equations(combo)
        prior_records.append(record(tier, cost, combo, status))
        if status == "sat-rational":
            repair = make_added_fields(combo, charges)
            assert verify_all_anomalies_vanish(repair)
            # Minimal only if every prior rung is support-pruned or unsat.
            return repair, certify_minimality(prior_records)
\end{lstlisting}

\section{Validation}
\label{sec:validation}

Because the anomaly/repair engine's value is its correctness, it was checked two
ways. First, the public anomaly/repair test suite (176 tests, distributed with
the code) pins the physics
against hand-computed values, including \emph{nonzero} anomaly coefficients, so
that a wrong Dynkin or multiplicity factor which still cancels for the symmetric
Standard Model cannot hide. Second, as internal quality assurance, the anomaly
and repair outputs were independently re-derived and differential-tested against
the implementation over randomized inputs; this process found and fixed real
defects (for example, a non-terminating uniqueness search and a crash on
algebraic solver models) and indicated that the full two-$U(1)$ check agrees
with a single-$U(1)$ check wherever the latter applies while additionally
catching mixed-anomaly failures, with no disagreement observed over
$2\times10^{4}$ random charge assignments. This internal cross-checking is
evidence of correctness, not a substitute for independent external reproduction;
the test suite is the reproducible artifact.

\paragraph{Known-result checks.}
As a validation of the same constraint machinery, \nsm{} can leave the Standard
Model hypercharges unknown, impose anomaly cancellation, Yukawa gauge
invariance, and the convention $Y_H=1$, and recover the unique one-generation
solution
\[
Y_{Q_L}=\tfrac13,\quad Y_{u_R}=\tfrac43,\quad Y_{d_R}=-\tfrac23,\quad
Y_{L_L}=-1,\quad Y_{e_R}=-2 .
\]
To isolate anomaly cancellation alone, the validation removes the Higgs and Yukawa
constraints and instead fixes $Y_{Q_L}=1/3$; the remaining solutions are exactly
the two assignments related by $u_R\!\leftrightarrow d_R$. Adding a free
right-handed neutrino opens the familiar $B\!-\!L$ flat direction, while a lone
colored Weyl fermion or an odd number of weak doublets is rejected by the
non-abelian cubic anomaly or Witten parity check, respectively.

\paragraph{Guaranteed versus observed.}
The engine rests on three trust anchors: anomaly coefficients for concrete
spectra are computed in exact rational arithmetic; \texttt{support-prune} lower bounds
are representation-theoretic subset checks; and the nontrivial solver claim
needed for minimality is a sound \texttt{unsat} result from Z3 for the asserted
real polynomial constraints~\cite{z3,nlsat}. A reported zero or nonzero
coefficient is therefore exact, and an SMT \texttt{unsat} verdict rules out the
tested repair combination over the reals. A lower-cost record made
only of \texttt{support-prune} and \texttt{unsat} verdicts is consequently a
solver-backed guarantee of minimality within the declared menu and cost
ordering. Everything else is reported as evidence, not proof: an
\texttt{unknown} or \texttt{sat-algebraic-model} rung does not certify
minimality, and differential agreement is a test result rather than a theorem.
For the anomaly and repair layers the tool either returns a guarantee or labels
the result uncertified. Optional downstream graph-search routines are
distributed with the package, but they are not part of the anomaly-repair
guarantee studied here.

\section{Related work}

\nsm{} is closest in spirit to model-building infrastructure, but it is aimed at
an earlier phase of the workflow. FeynRules~\cite{feynrules2} and
SARAH~\cite{sarah4} take a specified model (fields, symmetries, parameters, and
often a Lagrangian) and derive model data such as vertices, mass matrices,
renormalization-group equations, and interfaces to simulation tools. PyR@TE
focuses on automatic renormalization-group equations for general gauge
theories~\cite{pyrate}. SModelS decomposes BSM collider signatures into
simplified-model topologies and confronts them with LHC limits~\cite{smodels}.
These tools are deeper calculators and phenomenology interfaces than \nsm{} is.
A recent LLM-assisted BSM framework similarly keeps a deterministic symbolic
backend for Lagrangian construction, anomaly checks, operator expansion, and mass
matrices, with the LLM used as an orchestration layer~\cite{bsmagent}; it is
adjacent to, not a replacement for, the narrower inverse-repair workflow here.

The anomaly-equation literature addresses much closer mathematics.
Costa--Dobrescu--Fox give a general solution to the $U(1)$ anomaly
equations~\cite{u1solution}. Allanach--Davighi--Melville's
\emph{Anomaly-free Atlas} scans bounded $U(1)$ charge assignments and publishes
anomaly-free charge lists~\cite{anomalyatlas}. Allanach--Gripaios--Tooby-Smith
solve local anomaly equations for gauge-rank extensions with fixed SM-singlet
content~\cite{allanachsinglets}. Batra--Dobrescu--Spivak prove constructive
anomaly-free embeddings for arbitrary rational charges, including additional
matter that is vector-like under the SM but chiral under a new $U(1)$
factor~\cite{batra}. Rathsman--Tellander use algebraic-geometry methods to find
rational anomaly-free Froggatt--Nielsen-type models~\cite{rathsman}.

The distinction is the optimization problem and its certification record, both
relative to a declared SM representation menu. Prior work includes both
classification of anomaly-free assignments and constructive repair of anomalous
rational charge sets. Here the input is a fixed anomalous assignment, while the
output is the first exact rational repair found in a lexicographic menu ordering,
accompanied by support lower bounds, a record of every cheaper menu combination's
status, and an exact re-verification of the returned spectrum. The full
mixed-anomaly computation is essential infrastructure for that workflow, but not
presented as the principal novelty by itself.

\section{Limitations}
\label{sec:limits}

\nsm{} is a structural anomaly-repair engine, not a complete model builder. The
present paper treats a one-generation charge skeleton. When interpreted as a
three-family model, this corresponds to family-universal charges and does not
exploit possible intergenerational anomaly cancellation. The repair
menu is fixed to $N,E',L',D',Q'$ and the minimality claim is relative to that
menu and the representation-first cost ordering of Section~4. Charges are
reported as exact rationals; the SMT search over real polynomial constraints is
a relaxation used to obtain sound \texttt{unsat} results, and algebraic models
are marked uncertified for rational-repair purposes. The objective does not
optimize numerator or denominator complexity of the repaired charges.

The anomaly check covers perturbative local anomalies for the specified product
group, plus the modeled Witten $\SUtwo$ parity check. It does not cover general
discrete or global anomalies, nor Green--Schwarz or Wess--Zumino cancellation
mechanisms. The returned matter repair is only an anomaly-cancelling fermion
spectrum: \nsm{} does not synthesize the scalar sector, preserve or redesign
Yukawa interactions, generate masses for $\UX$-chiral exotics, compute cross
sections or decay rates, evaluate amplitudes, simulate detectors, or make
discovery claims. The optional process and graph-search modules distributed with
the package are downstream structural checks over a configured vertex subset;
they are not part of the paper's central guarantee.

\section{Conclusion}

\nsm{} demonstrates that anomaly diagnosis and minimal matter repair for
$\UX$ extensions can be turned into a reproducible structural workflow: find the
obstruction, test which repair menus can reach it, return an exact witness when
one is found, and say plainly where the evidence stops. The support bounds make
many repair-tier statements exact before the solver runs; the solver is reserved
for the remaining nonlinear repair search, and unresolved cases are labeled
uncertified. Natural extensions include emitting tracked unsat cores, SMT-LIB
replay artifacts, or solver proof logs for the minimality records, and, as a
separate project phase, moving from anomaly repair to full-model synthesis with
scalar sectors, Yukawa preservation, and chiral process labels.

\paragraph{Availability.} The implementation, tests, and census-generation
scripts are available at
\url{https://github.com/Morritse/nsm-anomaly-repair}, version
\texttt{0.1.0}, under the MIT License. The release corresponding to this
submission is tagged \texttt{paper1-cpc-submission}. Table~\ref{tab:census}
was generated with:
\begin{lstlisting}[language=bash]
.venv/bin/python scripts/repair_map.py \
    --radius 2 --max-added 6 --timeout-ms 8000 \
    --output csv --workers 16
\end{lstlisting}
using Python~3.14.4 and Z3~4.16.0 on Linux
\texttt{7.0.1-070001-generic x86\_64}. The generated census CSV has SHA-256:
\begin{lstlisting}
eca2f579baef8d91f11fcbe79273b6e907dd13d001136134c7487e33f9f4f244
\end{lstlisting}
The randomized differential checks described in Section~\ref{sec:validation} were
internal quality-assurance runs; the distributed reproducible artifacts are the
deterministic test suite and the repair-census command above. The release tag
and repository URL identify the software snapshot used here.

\section*{Declaration of generative AI and AI-assisted technologies in the manuscript preparation process}

During the preparation of this work, the author used
\href{https://github.com/Morritse/Claudex}{Claudex}, an open-source,
author-developed multi-agent orchestration interface, to coordinate Claude Code
and Codex sessions backed by Anthropic Claude and OpenAI models. These tools
were used to assist with drafting and language editing, structural organization
of the manuscript, and iterative critical review and cross-checking of the
manuscript's exposition, internal consistency, and framing of claims. After
using these tools, the author reviewed and edited the content as needed and
takes full responsibility for the content of the published article. The Claudex
source repository is available at \url{https://github.com/Morritse/Claudex}.

\begin{thebibliography}{99}
\small
\bibitem{z3} L.~de Moura and N.~Bj{\o}rner, \emph{Z3: An Efficient SMT Solver},
in TACAS 2008, LNCS \textbf{4963} (Springer), pp.~337--340 (2008).
\bibitem{nlsat} D.~Jovanovi{\'c} and L.~de Moura, \emph{Solving Non-Linear
Arithmetic}, in IJCAR 2012, LNCS \textbf{7364} (Springer), pp.~339--354 (2012).
\bibitem{feynrules2} A.~Alloul, N.~D.~Christensen, C.~Degrande, C.~Duhr, and
B.~Fuks, \emph{FeynRules 2.0--A complete toolbox for tree-level
phenomenology}, Comput.\ Phys.\ Commun.\ \textbf{185}, 2250--2300 (2014).
\bibitem{sarah4} F.~Staub, \emph{SARAH 4: A tool for (not only SUSY) model
builders}, Comput.\ Phys.\ Commun.\ \textbf{185}, 1773--1790 (2014).
\bibitem{pyrate} F.~Lyonnet, I.~Schienbein, F.~Staub, and A.~Wingerter,
\emph{PyR@TE: Renormalization Group Equations for General Gauge Theories},
Comput.\ Phys.\ Commun.\ \textbf{185}, 1130--1152 (2014).
\bibitem{smodels} S.~Kraml, S.~Kulkarni, U.~Laa, A.~Lessa, W.~Magerl,
D.~Proschofsky-Spindler, and W.~Waltenberger, \emph{SModelS: a tool for
interpreting simplified-model results from the LHC and its application to
supersymmetry}, Eur.\ Phys.\ J.\ C \textbf{74}, 2868 (2014).
\bibitem{u1solution} D.~B.~Costa, B.~A.~Dobrescu, and P.~J.~Fox,
\emph{General Solution to the $U(1)$ Anomaly Equations}, Phys.\ Rev.\ Lett.\
\textbf{123}, 151601 (2019).
\bibitem{anomalyatlas} B.~C.~Allanach, J.~Davighi, and S.~Melville,
\emph{An Anomaly-free Atlas: charting the space of flavour-dependent gauged
$U(1)$ extensions of the Standard Model}, JHEP \textbf{02}, 082 (2019),
arXiv:1812.04602.
\bibitem{allanachsinglets} B.~C.~Allanach, B.~Gripaios, and J.~Tooby-Smith,
\emph{Solving local anomaly equations in gauge-rank extensions of the Standard
Model}, Phys.\ Rev.\ D \textbf{101}, 075015 (2020), arXiv:1912.10022.
\bibitem{batra} P.~Batra, B.~A.~Dobrescu, and D.~Spivak,
\emph{Anomaly-Free Sets of Fermions}, J.\ Math.\ Phys.\ \textbf{47}, 082301
(2006), arXiv:hep-ph/0510181.
\bibitem{rathsman} J.~Rathsman and F.~Tellander,
\emph{Anomaly-free Model Building with Algebraic Geometry}, Phys.\ Rev.\ D
\textbf{100}, 055032 (2019), arXiv:1902.08529.
\bibitem{bsmagent} S.~Saad,
\emph{Large Language Model-Assisted Framework for BSM Model Building},
arXiv:2606.21316 (2026).
\bibitem{witten} E.~Witten, \emph{An $SU(2)$ Anomaly}, Phys.\ Lett.\ B
\textbf{117}, 324 (1982).
\bibitem{bminusL} R.~N.~Mohapatra and R.~E.~Marshak, \emph{Local $B\!-\!L$
Symmetry of Electroweak Interactions, Majorana Neutrinos and Neutron
Oscillations}, Phys.\ Rev.\ Lett.\ \textbf{44}, 1316 (1980).
\end{thebibliography}

\end{document}
